\section{Reconstrucción analítica de la energía y memoria del refinamiento}
\label{sec:limite-energetico-k}

Las separaciones positivas construidas a partir del registro dodecafásico
determinan una partición ordenada y, por la sección anterior, una forma
de Dirichlet sobre sus valores nodales. Su refinamiento posee dos
operaciones complementarias: la interpolación armónica conserva la energía
de los datos heredados, mientras cada detalle registra la variación
incorporada entre esos datos. El paso a profundidad ilimitada consiste en
reunir esta colección compatible mediante una norma que controle la suma de
sus detalles. El resultado concierne a la realización analítica
unidimensional de la energía; conserva como antecedente la construcción
HMT del registro y de la estructura discreta conjunta del continuo.

\subsection{Particiones anidadas y espacios de interpolación}

Fijemos la partición inicial
\(P_0=\{0=t_0<t_1<\cdots<t_{12}=1\}\), cuyas longitudes son
\(d_j=t_j-t_{j-1}=K_j/\sum_iK_i>0\). Sea
\((P_n)_{n\geq0}\) una sucesión de particiones finitas anidadas de
\([0,1]\), obtenidas por subdivisión de los intervalos anteriores.
Escribiremos
\[
 |P_n|=\max\{b-a:[a,b]\text{ es un intervalo de }P_n\},
 \qquad \lim_{n\to\infty}|P_n|=0.
\]
La subdivisión reiterada de cada intervalo en un número fijo de sucesores
iguales, al menos dos, satisface esta condición. Para otras reglas, la
contracción de la malla es una hipótesis explícita de este lector.
Las etiquetas de los sucesores y sus datos de acarreo acompañan la
subdivisión y conservan la procedencia del refinamiento.

Sea \(V_n\) el espacio complejo de funciones continuas sobre \([0,1]\)
que son afines en cada intervalo de \(P_n\). La energía y la norma
energética se definen por
\begin{equation}
 \mathcal E_n(f)
 =\sum_{[a,b]\in P_n}\frac{|f(b)-f(a)|^2}{b-a}
 =\int_0^1|f'(x)|^2\,dx,
 \qquad
 \|f\|_{\mathcal E}^2=|f(0)|^2+\mathcal E_n(f).
 \label{eq:energia-afines-k}
\end{equation}
La igualdad integral resulta de que la derivada de una función afín es
constante en cada intervalo. La energía coincide así exactamente con
la forma nodal ya construida. La inclusión \(V_n\subset V_{n+1}\)
interpreta la misma función afín en una partición más fina; conserva
tanto \(f(0)\) como la energía y es isométrica para
\(\|\cdot\|_{\mathcal E}\).

El valor inicial tiene una función precisa: fija la componente constante
que la forma de Dirichlet identifica en su núcleo. La energía controla
las diferencias, y el par formado por el valor inicial y esas diferencias
determina el observable. En el espacio de Sobolev
\(H^1([0,1];\mathbb C)\), cuyos elementos tienen representante
absolutamente continuo con derivada débil en \(L^2\), esta norma es
equivalente a la norma usual. En efecto,
\(f(x)=f(0)+\int_0^xf'(t)\,dt\) controla la norma de \(f\) a
partir de \(|f(0)|\) y \(\|f'\|_{L^2}\); la integración de la
misma identidad controla a su vez \(|f(0)|\) mediante
\(\|f\|_{L^2}+\|f'\|_{L^2}\).

\subsection{Reconstrucción de una función a partir de sus valores compatibles}

\begin{theorem}[Reconstrucción energética y completación de las interpolaciones]
\label{thm:limite-energia-h1-k}
Sea \(f_n\in V_n\) una colección compatible por restricción a los
nodos heredados:
\[
 f_{n+1}(a)=f_n(a)
 \quad\text{para todo nodo }a\text{ de }P_n.
\]
Si \(\sup_n\mathcal E_n(f_n)<\infty\), existe una única función
\(f\in H^1([0,1];\mathbb C)\), considerada en su representante
continuo, cuyos valores en todos esos nodos son los prescritos. Además,
\(f_n\to f\) uniformemente y en \(H^1\), y
\begin{equation}
 \int_0^1|f'(x)|^2\,dx
 =\lim_{n\to\infty}\mathcal E_n(f_n).
 \label{eq:limite-energia-h1-k}
\end{equation}
Recíprocamente, la interpolación nodal de cualquier elemento de \(H^1\)
produce una colección de esta clase. Por tanto, la completación de
\(\bigcup_n V_n\) en la norma energética es \(H^1([0,1];\mathbb C)\).
\end{theorem}

\begin{proof}
Pongamos \(g_n=f_n'\in L^2([0,1];\mathbb C)\). Para
\(m\geq n\), la compatibilidad de los valores en los extremos de cada
intervalo \([a,b]\) de \(P_n\) proporciona
\[
 \frac1{b-a}\int_a^bg_m(x)\,dx
 =\frac{f_m(b)-f_m(a)}{b-a}
 =\frac{f_n(b)-f_n(a)}{b-a}
 =g_n\big|_{(a,b)}.
\]
Sea \(W_n\) el espacio de funciones de \(L^2\) constantes en cada
intervalo de \(P_n\), y sea \(Q_n:L^2\to W_n\) la proyección
ortogonal dada por los promedios sobre dichos intervalos. La identidad
anterior afirma \(Q_ng_m=g_n\). En consecuencia,
\(g_m-g_n\) es ortogonal a \(g_n\) y
\begin{equation}
 \|g_m-g_n\|_{L^2}^2
 =\mathcal E_m(f_m)-\mathcal E_n(f_n).
 \label{eq:pitagoras-gradientes-k}
\end{equation}
Las energías son crecientes y acotadas. La identidad de Pitágoras
prueba que \((g_n)\) es una sucesión de Cauchy en \(L^2\), con
límite \(g\). Definamos
\[
 f(x)=f_0(0)+\int_0^xg(t)\,dt.
\]
Esta función pertenece a \(H^1\), tiene derivada débil \(g\) y
satisface, por la desigualdad de Cauchy--Schwarz,
\begin{equation}
 \sup_{x\in[0,1]}|f_n(x)-f(x)|
 \leq\|g_n-g\|_{L^2}\longrightarrow0.
 \label{eq:convergencia-uniforme-energia-k}
\end{equation}
La convergencia uniforme conserva cada valor nodal prescrito; junto con
la convergencia de las derivadas en \(L^2\), da la convergencia en
\(H^1\). La convergencia de las normas de las derivadas prueba
\eqref{eq:limite-energia-h1-k}. La unión de los nodos es densa porque
\(|P_n|\to0\); dos representantes continuos con los mismos valores
en esa unión coinciden. Se obtiene la unicidad.

Para la recíproca, tomemos \(f\in H^1\) en su representante
continuo y sea \(f_n\) su interpolación en los nodos de \(P_n\).
Su derivada es \(g_n=Q_nf'\). Las proyecciones \(Q_n\) son
contractivas en \(L^2\). Para una función continua \(v\), el
error \(|Q_nv-v|\) está acotado por su módulo de continuidad
evaluado en \(|P_n|\), que tiende a cero. La densidad de las
funciones continuas en \(L^2\) y la contractividad implican
\(Q_ng\to g\) para todo \(g\in L^2\): dado un aproximante
continuo \(v\), se emplea
\[
 \|Q_ng-g\|_{L^2}
 \leq 2\|g-v\|_{L^2}+\|Q_nv-v\|_{L^2}.
\]
Aplicando esta convergencia a \(g=f'\), se obtienen la convergencia
de las derivadas y, mediante \eqref{eq:convergencia-uniforme-energia-k},
las de las funciones. La contractividad acota sus energías por
\(\|f'\|_{L^2}^2\), y la convergencia de las normas prueba de
nuevo la identidad del límite. Toda función de \(H^1\) es, por
tanto, límite energético de elementos de \(\bigcup_nV_n\).
La equivalencia de normas y la completitud de \(H^1\) identifican
la completación enunciada.
\end{proof}

El argumento describe concretamente qué se conserva al pasar de la red
a la función. Cada derivada finita registra el promedio de todas las
derivadas posteriores en su intervalo; al refinar, se incorporan
componentes ortogonales que mantienen ese promedio. La energía acotada
controla la suma total de esas incorporaciones. El valor continuo se
recupera entonces integrando la derivada límite y conservando el valor
inicial. La realización analítica reúne, por esta vía explícita, la
información compatible de los niveles finitos.

\subsection{Descomposición ortogonal de la memoria por escalas}

Sea \(H_n\) la interpolación armónica de los valores de \(P_n\)
sobre los nodos de \(P_{n+1}\), interpretada como función de
\(V_{n+1}\), y definamos el detalle
\[
 z_{n+1}=f_{n+1}-H_nf_n.
\]
Por \eqref{eq:detalle-energia-k}, el detalle se anula en los nodos
heredados. Su derivada tiene promedio cero en cada intervalo antiguo
y es ortogonal a \(W_n\). Para niveles distintos, la derivada de
un detalle anterior pertenece a ese espacio de funciones constantes;
las derivadas de los detalles son, por tanto, ortogonales entre sí.
La misma ortogonalidad vale frente a \(f_0'\).

\begin{proposition}[Identidad de energía y recuperación de los detalles]
\label{prop:memoria-energetica-escalas-k}
Para toda profundidad finita \(N\),
\begin{equation}
 \mathcal E_N(f_N)
 =\mathcal E_0(f_0)
 +\sum_{n=0}^{N-1}\mathcal E_{n+1}(z_{n+1}).
 \label{eq:energia-telescopica-detalles-k}
\end{equation}
Si la colección satisface las hipótesis del
Teorema~\ref{thm:limite-energia-h1-k}, la serie de la derecha
converge y su suma es \(\int_0^1|f'|^2\). Los datos \(f_0\)
y todos los detalles determinan la colección compatible completa y su
realización analítica.
\end{proposition}
\begin{proof}
La identidad finita resulta de sumar
\(\mathcal E_{n+1}(f_{n+1})
=\mathcal E_n(f_n)+\mathcal E_{n+1}(z_{n+1})\), que es la
descomposición ortogonal de la sección anterior. Los sumandos son
no negativos y las sumas parciales convergen por
\eqref{eq:limite-energia-h1-k}. La reconstrucción finita se efectúa
recursivamente mediante \(f_{n+1}=H_nf_n+z_{n+1}\); el teorema
precedente determina después la función límite. Recíprocamente, una
función de \(H^1\) determina sus interpolaciones y, por diferencia,
todos estos detalles. Cada flecha conserva así sus datos propios.
\end{proof}

La eliminación variacional de nodos interiores selecciona el detalle
nulo en el nivel eliminado. La conservación de los detalles permite
recuperar cualquier configuración compatible de energía finita. La
respuesta efectiva y la configuración completa poseen, de esta manera,
una relación precisa: la primera registra el mínimo sobre las variables
interiores; la segunda añade las componentes ortogonales y su energía.

\subsection{Transporte de la realización analítica bajo el flujo positivo}

La dirección \(u=P_3P_{11}K\), recuperada antes de la evaluación
energética \cite{hmtX}, determina la colección de separaciones de
\eqref{eq:schur-flujo-k}. Su parámetro \(s\in\mathbb R\)
describe esta deformación analítica. Escribamos
\[
 d_j(s)=\frac{K_je^{su_j}}{\sum_\ell K_\ell e^{su_\ell}},
 \qquad
 t_i(s)=\sum_{j=1}^id_j(s),
 \qquad
 a_j(s)=\frac{d_j(s)}{d_j(0)}>0.
\]
Sea \(F_s:[0,1]\to[0,1]\) la aplicación que envía cada
\(t_i(0)\) a \(t_i(s)\) y es afín en los intervalos iniciales.
Es un homeomorfismo creciente, de pendiente \(a_j(s)\) en el
intervalo \([t_{j-1}(0),t_j(0)]\). Su inversa es también afín
por intervalos y ambas aplicaciones son Lipschitz.

Supongamos que cada subdivisión hereda las fracciones de longitud y la
dirección de su intervalo predecesor, como en la sección anterior. Todos
los nodos refinados se transportan entonces mediante el mismo \(F_s\).
Denotemos sus particiones por \(P_n(s)=F_s(P_n(0))\) y sus
espacios de interpolación por \(V_n(s)\). La aplicación
\[
 U_s f=f\circ F_s^{-1}
\]
transporta \(V_n(0)\) a \(V_n(s)\), preserva los valores
nodales correspondientes y prolonga esta acción a \(H^1\).

\begin{theorem}[Equivalencia energética y conmutación con el refinamiento]
\label{thm:transporte-h1-flujo-k}
Para todo \(f\in H^1([0,1];\mathbb C)\),
\begin{equation}
 \int_0^1|(U_sf)'(y)|^2\,dy
 =\sum_{j=1}^{12}\frac1{a_j(s)}
 \int_{t_{j-1}(0)}^{t_j(0)}|f'(x)|^2\,dx.
 \label{eq:transporte-energia-h1-k}
\end{equation}
Para \(s\) en cualquier intervalo compacto, las normas energéticas
antes y después del transporte son uniformemente equivalentes y
\(|P_n(s)|\to0\) uniformemente en ese compacto. La interpolación
armónica y la restricción a nodos heredados conmutan con \(U_s\);
la reconstrucción y la minimización de los detalles se transportan
tanto a profundidad finita como en la realización \(H^1\).
\end{theorem}

\begin{proof}
En el intervalo inicial \(j\), la regla de la cadena y el cambio
de variable \(y=F_s(x)\) proporcionan
\((U_sf)'(F_s(x))=f'(x)/a_j(s)\) y \(dy=a_j(s)\,dx\).
La integración y la suma sobre los intervalos prueban
\eqref{eq:transporte-energia-h1-k}. La fórmula vale para los
representantes absolutamente continuos, y su lado derecho es finito.

Fijado un intervalo compacto de parámetros, la continuidad y positividad
de los doce factores dan constantes \(0<a_-\leq a_+<\infty\)
tales que \(a_-\leq a_j(s)\leq a_+\) para todo \(j\) y
todo \(s\) de ese compacto. Por tanto,
\[
 \frac1{a_+}\mathcal E(f)
 \leq\mathcal E(U_sf)
 \leq\frac1{a_-}\mathcal E(f),
 \qquad
 |P_n(s)|\leq a_+|P_n(0)|.
\]
Como \((U_sf)(0)=f(0)\), se obtiene además
\[
 \min\{1,a_+^{-1}\}\|f\|_{\mathcal E}^2
 \leq\|U_sf\|_{\mathcal E}^2
 \leq\max\{1,a_-^{-1}\}\|f\|_{\mathcal E}^2.
\]
Estas cotas prueban la equivalencia uniforme y la conservación de la
condición de malla decreciente.

En cada arista predecesor, los coeficientes de interpolación son las sumas
parciales de las fracciones heredadas, independientes de \(s\).
Si \(H_n(s)\) denota su interpolación, se tiene
\(U_sH_n(0)=H_n(s)U_s\) en los espacios correspondientes.
La restricción conmuta porque los nodos se transportan con sus valores.
La descomposición en función armónica y detalle se conserva, y la
identidad de Schur de \eqref{eq:schur-flujo-k} preserva su
caracterización variacional. Finalmente, la continuidad acotada de
\(U_s\) permite pasar al límite \(H^1\). En particular, una
colección reconstruida antes del transporte converge después a \(U_sf\),
con los mismos valores nodales transportados.
\end{proof}

La ley \eqref{eq:transporte-energia-h1-k} identifica los factores
métricos exactos de la deformación: el cambio de longitud de cada
intervalo modifica inversamente su contribución energética. La
equivalencia de normas conserva las configuraciones de energía finita,
y el archivo de detalles conserva su reconstrucción. Se obtiene así una
realización analítica completa de este lector del registro positivo.
Su compatibilidad con las formas canónicas se articula sobre las
subdivisiones de un dominio fijo: la suma de formas conserva el objeto
meromorfo, mientras la reducción variacional y los detalles conservan
la respuesta energética y el observable refinado. La identificación
con formas amplituhedrales de rango superior conserva las condiciones
específicas de cobertura, grado y correspondencia de residuos
establecidas para aquella realización.
